\section*{Chapter \ref{ch:b2:climate-proxies} --- Climate Archives and Proxies}

\begin{solution}{exo:b2:climate-proxies:1}
(a) $\delta = (27.5 \times (-30) + 1350 \times 0)/(1350 + 27.5) = -825/1377.5 = -0.60$ per mil.
(b) The ocean would then hold all the water: $-0.60$ per mil, \num{0.60} per mil lighter than
today. (c) $27.5/1377.5 = \qty{2.0}{\percent}$.
\end{solution}

\begin{solution}{exo:b2:climate-proxies:2}
(a) $T = 16.1 - 4.76\,(2.50 - (0 - 0.27)) = 16.1 - 4.76 \times 2.77 = 16.1 - 13.19 =
\qty{2.9}{\celsius}$. (b) $T = 16.1 - 4.76\,(2.50 - 0.73) = 16.1 - 8.43 = \qty{7.7}{\celsius}$: the
same shell gives a temperature almost five degrees higher if the water was heavier, which is why
the composition of the water must be known. (c) The shell is $2.71 - 2.50 = 0.21$ per mil heavier,
so its water was $4.76 \times 0.21 = \qty{1.0}{\celsius}$ colder, at \qty{1.9}{\celsius}.
\end{solution}

\begin{solution}{exo:b2:climate-proxies:3}
(a) $0.4\,\mathrm{e}^{0.9} = \qty{0.98}{mmol/mol}$ and $0.4\,\mathrm{e}^{1.8} =
\qty{2.42}{mmol/mol}$. (b) The factor is $\mathrm{e}^{0.9} = 2.46$; per degree $\mathrm{e}^{0.09}
= 1.094$, a growth of \qty{9.4}{\percent} per degree. (c) $T = \ln(\text{Mg/Ca}/0.4)/0.09$:
$\ln 5/0.09 = \qty{17.9}{\celsius}$ and $\ln 2.5/0.09 = \qty{10.2}{\celsius}$.
\end{solution}

\begin{solution}{exo:b2:climate-proxies:4}
(a) $2668 - 2213 + 1 = 456$ years. (b) A single tree may lack a ring locally, or show a ring
that cannot be read; matching the patterns of wide and narrow rings across many trees reveals the
missing rings and fixes the year of every ring, and the agreement between trees measures the
reliability of the dating. (c) At the tree line one factor, the cold of the growing season,
limits all the trees, so they share the same pattern of good and bad years; in a mild, moist
lowland no single factor dominates, the patterns of the trees differ, and \omterm{def:b2:climate-proxies:dendro}{cross-dating} may fail.
\end{solution}

\begin{solution}{exo:b2:climate-proxies:5}
(a) $125 \times 40/(3682.2 - 125) = 5000/3557.2 = 1.41$ per mil; $134 \times 40/3548.2 =
5360/3548.2 = 1.51$ per mil. (b) With $H - h = 3552.2$: $3900/3552.2 = 1.10$, $4550/3552.2 = 1.28$
and $6500/3552.2 = 1.83$ per mil. (c) $\delta_{\mathrm{ice}} = -1.0 \times 3552.2/130 = -27.3$ per
mil, heavier than any present ice sheet (Greenland $-30$ to $-35$, the North American ice sheet
$-40$, Antarctic ice as low as $-50$). The mass balance with realistic ice and the pore-water
measurement do not agree; the gap of about half a per mil is a real uncertainty of the method.
\end{solution}

\begin{solution}{exo:b2:climate-proxies:6}
(a) $0.372\,\mathrm{e}^{0.0936T} = 0.4\,\mathrm{e}^{0.09T}$ gives $\mathrm{e}^{0.0036T} =
0.4/0.372 = 1.0753$, so $T = \ln 1.0753/0.0036 = \qty{20.2}{\celsius}$. (b) At
\qty{2}{\celsius}: $0.4\,\mathrm{e}^{0.18} = 0.479$ and $0.372\,\mathrm{e}^{0.1872} = 0.449$, so the
fitted ratio is $0.449/0.479 = 0.937$, \qty{6.3}{\percent} lower. (c) $T = \ln(0.449/0.4)/0.09 =
\ln 1.1225/0.09 = \qty{1.3}{\celsius}$: the predicted law reads the shell about
\qty{0.7}{\celsius} too cold. In cold water the choice of calibration matters at the level of a
degree.
\end{solution}

\begin{solution}{exo:b2:climate-proxies:7}
(a) From the Holocene to GI-1, \qty{112.19}{m} of ice for 2989 years: $11\,219/2989 =
\qty{3.75}{cm}$ per year; from GI-16 to GI-17, \qty{17.89}{m} for 1160 years: $1789/1160 =
\qty{1.54}{cm}$ per year. (b) The ice
flows: as each layer is buried it spreads sideways and thins, and the deeper, older layers have
been thinned longer. (c) When the layers become too thin, a centimetre or less, the instruments
can no longer resolve one year from the next, and the \omterm{def:b2:climate-proxies:age-model}{age model} must use tie points or flow models
instead.
\end{solution}

\begin{solution}{exo:b2:climate-proxies:8}
(a) The fraction of the interval is $(1550 - 1492.45)/(1604.64 - 1492.45) = 57.55/112.19 = 0.513$,
so the age is $11\,703 + 0.513 \times 2989 = 11\,703 + 1533 = \num{13236}$ years b2k. (b) 1950 is 50
years before 2000: \num{13186} years before 1950. (c) $50 + 0.513 \times (93 - 50) = 72$ years.
\end{solution}

\begin{solution}{exo:b2:climate-proxies:9}
(a) Each uncertain layer contributes half a year: $2 \times 93 = 186$ and $2 \times 1287 = 2574$
uncertain layers. (b) $\frac12\sqrt{186} = 6.8$ years and $\frac12\sqrt{2574} = 25.4$ years.
(c) $93/6.8 = 14$ and $1287/25.4 = 51$. The ratio of $\frac12 n$ to $\frac12\sqrt{n}$ is
$\sqrt{n}$: it grows as the square root of the number of uncertain layers, so the choice between
the two rules matters more and more with depth.
\end{solution}

\begin{solution}{exo:b2:climate-proxies:10}
(a) On the \omterm{prop:b2:isotope-geochemistry:mwl}{meteoric water line} $\delta$D changes eight times as much as $\delta^{18}$O, so the
slope is $8 \times 0.8 = 6.4$ per mil of $\delta$D per degree. The drop is $441.7 - 396.6 = 45.1$
per mil, hence a cooling of $45.1/6.4 = \qty{7.0}{\celsius}$. (b) The glacial ocean was about
$8 \times 1.0 = 8$ per mil heavier in $\delta$D; snow formed from today's ocean would have been
lighter by as much, $-441.7 - 8 = -449.7$ per mil. The drop becomes $53.1$ per mil and the cooling
$53.1/6.4 = \qty{8.3}{\celsius}$. (c) $9.39 - 0.09 = \qty{9.30}{\celsius}$, so $45.1/9.30 = 4.85$
per mil per degree. The published series is also corrected for the changing elevation of the ice
sheet, and its temperature relation comes from models that simulate the water isotopes, not from
the spatial slope alone.
\end{solution}

\begin{solution}{exo:b2:climate-proxies:11}
(a) The logarithms of the ratios are $-0.528$, $-0.223$, $0.191$, $0.501$, $0.904$, $1.227$; their
mean is $0.345$ and the mean temperature 14. With $S_{TT} = \sum (T - 14)^2 = 280$ and $S_{T\ell} =
\sum (T - 14)(\ell - 0.345) = 24.93$, $A = 24.93/280 = \qty{0.0890}{\per\celsius}$ and $\ln B =
0.345 - 0.0890 \times 14 = -0.901$, so $B = \qty{0.406}{mmol/mol}$. (b) The residuals are
$+0.017$, $-0.034$, $+0.023$, $-0.023$, $+0.025$ and $-0.009$, a few per cent of the ratio, with no
trend; the fit is close to $0.4\,\mathrm{e}^{0.09T}$. (c) $T = (\ln 1.30 + 0.901)/0.0890 = (0.262 +
0.901)/0.0890 = \qty{13.1}{\celsius}$. (d) $T = (\ln 0.45 + 0.901)/0.0890 = (-0.799 +
0.901)/0.0890 = \qty{1.2}{\celsius}$. This value lies below the coldest sample of the calibration,
\qty{4}{\celsius}: the law is extrapolated, and at low temperatures the Mg/Ca of shells is known to
depart from the exponential law under the influence of the carbonate ion.
\end{solution}

\begin{solution}{exo:b2:climate-proxies:12}
(a) Per unit of ocean area, the water of the present ocean, $H$ at 0, is shared between the stores
$h_i$ at $\delta_i$ and the remaining ocean $H - h$, with $h = \sum h_i$: $0 = \sum h_i\delta_i +
(H - h)\,\Delta\delta$. After the first stage $\Delta\delta = 60 \times 30/(3682.2 - 60) =
1800/3622.2 = 0.50$ per mil; after the second $(1800 + 70 \times 45)/(3682.2 - 130) =
4950/3552.2 = 1.39$ per mil. (b) $-4950/130 = -38.1$ per mil. (c) The rule is $1.39/130 = 0.0107$
per mil per metre, and gives $0.50/0.0107 = \qty{46}{m}$ for the first stage instead of
\qty{60}{m}: the first ice was heavier than the mean, so each of its metres shifted the ocean
less. A linear rule calibrated on one state misreads the others when the mean composition of the
ice changes.
\end{solution}

\begin{solution}{pb:b2:climate-proxies:1}
\textbf{1.} $361.84 \times 10^6\ \unit{km^2} \times \qty{0.130}{km} = \qty{4.70e7}{km^3}$, about 47
million cubic kilometres.
\textbf{2.} $130 \times 40/(3682.2 - 130) = 5200/3552.2 = 1.46$ per mil.
\textbf{3.} $5000/3557.2 = 1.41$ and $5360/3548.2 = 1.51$ per mil.
\textbf{4.} $4550/3552.2 = 1.28$ and $6500/3552.2 = 1.83$ per mil. The composition of the ice
moves the answer by \num{0.55} per mil, the sea level by \num{0.10}: the ice matters more.
\textbf{5.} $\delta_{\mathrm{ice}} = -1.0 \times 3552.2/130 = -27.3$ per mil.
\textbf{6.} $40/3682.2 = 0.0109$ per mil per metre, which rounds to the \num{0.011} of the dated
sea levels.
\textbf{7.} At constant $\delta^{18}\mathrm{O}_{\mathrm{sw}}$, $\mathrm{d}T = -4.76\,\mathrm{d}\delta^{18}
\mathrm{O}_{\mathrm{b}}$, so $\mathrm{d}T = -1$ gives $\mathrm{d}\delta^{18}\mathrm{O}_{\mathrm{b}} =
1/4.76 = 0.21$ per mil.
\textbf{8.} $\delta^{18}\mathrm{O}_{\mathrm{b}} = (16.1 - T)/4.76 + \delta^{18}\mathrm{O}_{\mathrm{sw}}
- 0.27 = 14.1/4.76 - 0.27 = 2.96 - 0.27 = 2.69$ per mil.
\textbf{9.} $17.1/4.76 + 1.0 - 0.27 = 3.59 + 0.73 = 4.32$ per mil. The change is $4.32 - 2.69 =
1.63$ per mil: \num{1.00} from the water and $3.0 \times 0.21 = 0.63$ from the cooling of
\qty{3}{\celsius}.
\textbf{10.} $T = \ln(2.0/0.4)/0.09 = \ln 5/0.09 = \qty{17.9}{\celsius}$.
\textbf{11.} $0.18/0.1 = \qty{1.8}{\celsius}$ of cooling.
\textbf{12.} Mg/Ca is a ratio of two elements in the calcite, set by the temperature of
calcification; it does not involve the isotopes of oxygen, and so it ignores the composition of
the water that the ice changes.
\textbf{13.} $4.92 - 3.23 = 1.69$ per mil.
\textbf{14.} Solving $\Delta = h \times 40/(H - h)$ for $h$: $h = H\Delta/(\Delta + 40) = 3682.2
\times 1.69/41.69 = \qty{149}{m}$, more than the 125 to \qty{134}{m} of the sea-level evidence:
part of the change must be cooling.
\textbf{15.} $1.69 - 1.46 = 0.23$ per mil, and $4.76 \times 0.23 = \qty{1.1}{\celsius}$.
\textbf{16.} $1.69 - 1.00 = 0.69$ per mil, and $4.76 \times 0.69 = \qty{3.3}{\celsius}$ (3.28).
\textbf{17.} $0.23/0.25 = \qty{0.9}{\celsius}$ and $0.69/0.25 = \qty{2.8}{\celsius}$.
\textbf{18.} The North Atlantic Deep Water cooled from 2--\qty{4}{\celsius} to \qty{-1.1}{\celsius},
by 3.1 to \qty{5.1}{\celsius}; the Antarctic Bottom Water from 0--\qty{1}{\celsius} to
\qty{-1.9}{\celsius}, by 1.9 to \qty{2.9}{\celsius}; the Mg/Ca of the North-West Atlantic gives
\qty{1.8}{\celsius}. All lie between about 2 and \qty{5}{\celsius}, around the 2.8 to
\qty{3.3}{\celsius} of the pore-water correction, and above the \qty{1}{\celsius} of the mass
balance with ice at $-40$ per mil.
\textbf{19.} $-396.6 - (-441.7) = 45.1$ per mil.
\textbf{20.} $8 \times 0.8 = 6.4$ per mil per degree, so $45.1/6.4 = \qty{7.0}{\celsius}$.
\textbf{21.} $9.39 - 0.09 = \qty{9.30}{\celsius}$, and $45.1/9.30 = 4.85$ per mil per degree.
\textbf{22.} Dome C, \qty{9.3}{\celsius}, more than the global surface, 5 to \qty{7}{\celsius}, more
than the deep ocean, about \qty{3}{\celsius}.
\textbf{23.} The deep water is already cold today, 0 to \qty{4}{\celsius}, and sea water cannot
cool much below its freezing point, which the glacial bottom water nearly reached at
\qty{-1.9}{\celsius}: the deep ocean had only a few degrees to lose, while the air over Antarctica
has no such floor.
\textbf{24.} With the pore-water correction, the deep ocean was about \textbf{\qty{3.3}{\celsius}
colder} at the Last Glacial Maximum than today, close to the 2 to \qty{5}{\celsius} of the
independent deep-water estimates.
\end{solution}
